% Original: PDF strana 25, gornji slajd; odobren drugi deo s049.
\slajdid{02-s049b}
\begin{frame}[label=02-s049b]{Dvoulazna NILI kola i složena CMOS realizacija}
\setlength{\parskip}{0pt}
% element: 02-s049b:e01
Troulazno NILI kolo zamenjujemo sa dva dvoulazna NILI kola i invertorom:
\par\vspace{2mm}
{\centering\(\displaystyle\overline{X+Y+Z}=\overline{(X+Y)+Z}
=\overline{\overline{\overline{X+Y}}+Z}.\)\par}
\par\vspace{3mm}
% element: 02-s049b:e02
Za realizaciju samo dvoulaznim NILI kolima potrebni su:
\begin{itemize}
\item 4 kola za ulazne inverzije — NILI sa spojenim ulazima;
\item 3 kola za zamenu troulaznog NILI, uključujući međuinverziju;
\item 3 kola za preostali deo funkcije.
\end{itemize}
Ukupno: \(4+3+3=10\) dvoulaznih NILI kola. Četiri kola po čipu: \textbf{3 ista čipa}.
\par\vspace{3mm}
% element: 02-s049b:e03
\Rezultat{\textbf{Probajte:} uz raspoložive \alert{\textbf{komplementne vrednosti ulaza}}, uporedite oba oblika funkcije kao \alert{\textbf{jednostepena složena CMOS kola}}. Koja realizacija zauzima \alert{\textbf{manju površinu}}?}
\end{frame}
